\chapter{المخرج الثاني: كيف يتحكم الدماغ في كيمياء الجسم}
\chaptermark{المخرج الكيميائي}
\label{sec:chemoutput}

\section{مخرج سريع ومخرج بطيء}

رأينا في الفصل~\ref{sec:muscles} المخرج السريع للآلة: العضلات. لكن للآلة مخرجًا ثانيًا بطيئًا. فالدماغ يُبقي حرارة الجسم ومعدل ضربات القلب وكمية الماء والسكر ضمن الحدود المطلوبة، ويهيّئ الجسم للركض أو للنوم. ولا يحرّك لذلك المفاصل، بل يطلق إشارات كيميائية. وللإشارات طريقان. الأول أعصاب تمتد مباشرة إلى الأعضاء الداخلية: إلى القلب والأوعية الدموية والأمعاء والغدد. والثاني الدم: بعض العصبونات والغدد لا تطلق مادتها في الشق نحو الجار، بل في مجرى الدم مباشرة.

الدم أوسع نواقل البث بين كل ما مرّ في هذا الكتاب. فقنوات البث في الفصول~\babelsublr{\ref{sec:serocomputer}--\ref{sec:acet}} كانت ترسل الإشارة إلى مناطق واسعة من الدماغ؛ أما الدم فيرسلها إلى كل خلية في الجسم. يدور الدم في الجسم في نحو دقيقة، فتصل الرسالة إلى الجميع خلال عشرات الثواني وتبقى دقائق وساعات حتى تتفكك المادة. وليس لهذا البث عنوان على الإطلاق. وكما في القضية~\ref{prop:receptor}، يختار المتلقّي معنى الرسالة: لا تستجيب إلا الخلايا التي تملك مستقبلًا لهذه المادة.

\section{المسرِّع والمكبح: سلكان إلى الأعضاء}

خير مثال القلب. في القلب ناظمة إيقاع خاصة به، كما في عصبون \nt{SERO} في الفصل~\ref{sec:serocomputer}: فهي تضع الإيقاع بنفسها، ولو قُطعت كل الأعصاب لنبض القلب نحو مئة مرة في الدقيقة. الدماغ لا يولّد الإيقاع، بل يضبطه فقط، وذلك عبر سلكين متعاكسي الإشارة (الشكل~\ref{fig:gasbrake}). في أحدهما يسري \nt{NORA}: وهو \emph{المسرِّع}، يرفع الإيقاع. وفي الآخر \nt{ACET}: وهو \emph{المكبح}، يخفضه. تقريبًا،
\begin{empheq}[box=\keybox]{equation}
f \;\approx\; f_0 \;+\; a_S\,S \;-\; a_P\,P ,
\label{eq:gasbrake}
\end{empheq}
حيث $f_0\approx100$ ضربة في الدقيقة هو الإيقاع الذاتي لناظمة الإيقاع، و$S$ و$P$ نشاط المسرِّع والمكبح. في الراحة يكون المكبح مضغوطًا، فينبض القلب عادةً نحو ستين إلى سبعين مرة في الدقيقة؛ وعند الجهد يُرفع المكبح ويُضغط المسرِّع.

\begin{figure}[t]
\centering
\begin{tikzpicture}[>=Stealth,
  blk/.style={draw,very thick,rounded corners=2pt,align=center,font=\small},
  pace/.style={circle,draw,very thick,double,double distance=1.2pt,minimum size=1.8cm,inner sep=0pt,font=\scriptsize,fill=red!8,align=center},
  inh/.style={-{Bar[width=7pt]},thick,red!70!black}]
\node[blk,fill=blue!5,text width=1.6cm] (B) at (0,0) {أمر\\الدماغ};
\node[pace] (H) at (6.2,0) {ناظمة\\الإيقاع};
\draw[->,very thick,orange!85!black] (B.north east) to[bend left=20] node[above,font=\scriptsize,black,align=center] {\nt{NORA}: مسرِّع، ثوانٍ} (H.north west);
\draw[inh,very thick,blue!60!black] (B.south east) to[bend right=20] node[below,font=\scriptsize,black,align=center] {\nt{ACET}: مكبح، أجزاء من الثانية} (H.south west);
\draw[->,very thick] (H) -- (8.4,0) node[right,font=\small] {التردد $f$};
\node[blk,fill=orange!10,text width=1.6cm] (A) at (0,-2.6) {\nt{ADRE}\\إلى الدم};
\draw[->,thick,orange!85!black] (B) -- (A);
\foreach \y/\lab in {-2.0/{القلب},-2.6/{الأوعية},-3.2/{الكبد، العضلات}}{
  \draw[->,thick,densely dashed,orange!85!black] (A.east) -- (4.3,\y) node[right,font=\scriptsize,black] {\lab};}
\end{tikzpicture}
\caption{سلكان متعاكسا الإشارة إلى ناظمة إيقاع القلب: المسرِّع \nt{NORA} بطيء، والمكبح \nt{ACET} سريع. في الأسفل المسرِّع نفسه عبر ناقل البث: تطلق خلايا \nt{ADRE} الأدرينالين في الدم، فتتلقى الإشارة كل الأعضاء التي تملك المستقبل المناسب (الأسهم المتقطعة).}
\label{fig:gasbrake}
\end{figure}

هذه هي الشفرة ثنائية السلك في الفصل~\ref{sec:glutcomputer} وزوج العضلات في الفصل~\ref{sec:muscles}، غير أن السلكين الآن لا يتحكمان في مفصل، بل في سرعة ناظمة الإيقاع. وللسلكين سرعتان مختلفتان. كلاهما يعمل مرشحَ تمرير منخفض، لكن المكبح يمرّر التغيرات أسرع من المسرِّع بعدة أضعاف~\cite{Berger1989}: طريق \nt{ACET} قصير، إذ يفتح البروتين المرتبط بالمستقبل قناة البوتاسيوم بنفسه، دون مرسال ثانٍ داخل الخلية~\cite{Logothetis1987gbg}، أما \nt{NORA} فيعمل عبر سلسلة من التفاعلات داخل الخلية. يعرف المهندس هنا حيلة المشغّلَين المختلفَي السرعة: السريع يتولى الضبط الدقيق الفوري، والبطيء يتولى التغيرات الكبيرة الطويلة.

وهنا يجد \nt{ADRE} دوره، وقد كُتب عنه في الجدول~\ref{tab:types} أن دوره في الدماغ غير واضح. أما خارج الدماغ فدوره واضح. بعض خلايا طريق المسرِّع لا ترسل سلكًا إلى عضو، بل تطلق الأدرينالين في الدم مباشرة. إنه المسرِّع نفسه، لكن عبر ناقل البث: يصل خلال عشرات الثواني إلى القلب والأوعية والكبد والعضلات معًا ويبقى دقائق. المسرِّع المعنوَن \nt{NORA} يضبط عضوًا واحدًا؛ والمسرِّع المبثوث \nt{ADRE} ينقل الجسم كله إلى وضع الركض.

\section{سلسلة مراحل بتغذية راجعة}

كثير من الهرمونات لا تطلقها غدة واحدة، بل سلسلة من المراحل. تطلق مجموعة صغيرة من العصبونات في قاعدة الدماغ الإشارة الأولى في الدم؛ فتدفع غدةً وسيطة إلى إطلاق إشارة ثانية؛ وهذه تدفع الغدة الأخيرة إلى إطلاق الهرمون الذي يؤثر في الجسم. ثم يعود الهرمون الأخير إلى المراحل الأولى فيثبطها. فينتج مضخم من ثلاث مراحل تحيط به تغذية راجعة سالبة، أي منظِّم مستوى، كما في منظومة \nt{SERO} في الصيغة~\eqref{eq:feedback}.

لنصف كل مرحلة بدارة RC ثابتها الزمني $\tau$ وكسبها $g$، والتغذية الراجعة بالمعامل $k$:
\begin{equation}
\tau\,\frac{\dd x_1}{\dd t} = -x_1 + g\bigl[u - k\,x_3\bigr]_+ ,\qquad
\tau\,\frac{\dd x_2}{\dd t} = -x_2 + g\,x_1 ,\qquad
\tau\,\frac{\dd x_3}{\dd t} = -x_3 + g\,x_2 ,
\label{eq:cascade}
\end{equation}
حيث $u$ أمر الدماغ، و$x_3$ مستوى الهرمون الأخير. كسب الحلقة $K=g^3k$. في التوازن $x_3=g^3u/(1+K)$، وكما في أي منظِّم بتغذية راجعة، تجعل الحلقةُ الكبيرة المستوى غير حساس لتفاوت كسب المراحل. لكن كل مرحلة تزيح الطور: عند التردد $\omega=\sqrt3/\tau$ تعطي المراحل الثلاث المتماثلة معًا إزاحة نصف دورة وتوهينًا بثمانية أضعاف. ومن هنا النتيجة الكلاسيكية في نظرية التحكم:
\begin{empheq}[box=\keybox]{equation}
K \;<\; 8 ,
\label{eq:cascadestable}
\end{empheq}
وإلا بدأ المنظِّم يتذبذب بدور يقارب $2\pi\tau/\sqrt3\approx3.6\,\tau$.

\begin{keyblock}
\begin{proposition}[الدقة مقابل الاستقرار]
\label{prop:cascade}
في سلسلة من ثلاث مراحل بتغذية راجعة تناسبية، خطأ المستوى يساوي $1/(1+K)$، ولا تستقر السلسلة إلا عندما $K<8$. لذلك لا يستطيع هذا المنظِّم أن يثبّت المستوى بدقة أفضل من نحو عشرة في المئة؛ ومحاولة رفع الكسب تحوّله إلى مولّد تذبذبات.
\end{proposition}
\end{keyblock}

تحققنا من ذلك على النموذج~\eqref{eq:cascade} (الشكل~\ref{fig:cascade}). عند $K=4$ يتذبذب المستوى قليلًا بعد التشغيل ثم يستقر. وعند $K=7$ يستقر أيضًا، لكن ببطء. وعند $K=12$ لا تخمد التذبذبات ولا تنمو بلا حد: فالمرحلة لا تستطيع إخراج إشارة سالبة، فتدخل السلسلة في نبضات منتظمة بين $0.83$ و$1.24$ من المستوى المتوسط بدور $3.7\text{--}3.9\,\tau$.

\begin{figure}[t]
\centering
\begin{tikzpicture}[>=Stealth,x=0.3cm,y=1.2cm]
\draw[->,gray!70!black] (0,0) -- (41.5,0) node[right,font=\small,black] {$t/\tau$};
\draw[->,gray!70!black] (0,0) -- (0,2.8) node[left,font=\small,black] {$x_3/x_3^*$};
\foreach \t in {0,10,20,30,40}{\draw[gray!60!black] (\t,-0.05) -- (\t,0.05); \node[below,font=\scriptsize] at (\t,0) {\t};}
\foreach \f in {1,2}{\draw[gray!60!black] (-0.3,\f) -- (0.3,\f); \node[left,font=\scriptsize] at (0,\f) {\f};}
\draw[densely dashed,gray!60!black] (0,1) -- (40,1);
\draw[thick,blue!60!black] plot file {figures/cascade_K4.dat};
\draw[thick,red!70!black] plot file {figures/cascade_K12.dat};
\node[font=\scriptsize,blue!60!black,anchor=west] at (16,0.55) {$K=4$: يستقر};
\node[font=\scriptsize,red!70!black,anchor=west] at (16,1.75) {$K=12$: ينبض};
\end{tikzpicture}
\caption{سلسلة من ثلاث مراحل بتغذية راجعة وفق~\eqref{eq:cascade}؛ مستوى الهرمون الأخير بوصفه كسرًا من قيمة التوازن $x_3^*$. حين يكون كسب الحلقة أقل من ثمانية يستقر المستوى؛ وحين يزيد عليها تولّد السلسلة بنفسها نبضات منتظمة.}
\label{fig:cascade}
\end{figure}

والهرمونات الحقيقية في هذه السلاسل تخرج فعلًا على شكل نبضات: فمستوى هرمون الإجهاد مثلًا يتذبذب بدور يقارب الساعة. وبحسب إحدى الفرضيات، تولّد هذه النبضاتِ التغذيةُ الراجعة نفسها بين المراحل مع تأخرها، ولا حاجة إلى مولّد إيقاع مستقل~\cite{Walker2010}. وهذا هو سبب التذبذب نفسه في حلقة \foreignlanguage{english}{\nt{GLUT}--\nt{GABA}} في القضية~\ref{prop:ei} وفي حلقة شدّ العضلة في الشرط~\eqref{eq:delaystable}: تغذية راجعة سالبة متأخرة.

\section{الشفرة النبضية للهرمونات}

النبضات ليست أثرًا جانبيًا فحسب، بل قد تكون شفرة. العصبونات التي تتحكم في التكاثر تطلق إشارتها في الدم دفعاتٍ قصيرة نحو مرة كل ساعة. فإذا أُعطيت الغدة المتلقية الإشارةَ نفسها باستمرار لا دفعاتٍ، فإنها لا تعمل بكامل طاقتها كما مع النبضات، بل تصمت؛ وإذا عادت الإشارة دفعاتٍ مرة كل ساعة، عاد العمل~\cite{Belchetz1978}. إذن المتلقّي لا يقرأ المستوى، بل \emph{تردد} الدفعات.

لا لغز في ذلك عند المهندس. المستقبل الذي يتعب ويكفّ عن الاستجابة تحت تأثير المادة الطويل هو مرشح تمرير عالٍ، مثل المشبك المستنفَد~\eqref{eq:depression} في الفصل~\ref{sec:code}. فهو يخمد الإشارة المستمرة ويمرّر التغيرات. ولا يمكن إبلاغ مثل هذا المتلقّي شيئًا إلا بالنبضات، فتنتقل المعلومة من المستوى إلى تردد الدفعات وشكلها. ثم إن الترددات المختلفة للدفعات تؤثر تأثيرات مختلفة في خلايا مختلفة من الغدة نفسها، فيمكن نقل أكثر من كمية واحدة عبر خط كيميائي واحد، كما سرت مركّبتان سريعة وبطيئة عبر سلك \nt{DOPA} واحد في الفصل~\ref{sec:dofaalt}.

\section{الإيقاع اليومي: مولّد بطيء للأوضاع}

أبطأ إيقاعات الجسم هو الإيقاع اليومي. تولّده تغذية راجعة سالبة متأخرة داخل الخلايا: تنتج الجينات بروتينات تثبّط إنتاجها هي نفسها بتأخر عدة ساعات~\cite{TakahashiJS2017}. مرة أخرى تغذية راجعة سالبة متأخرة، للمرة الرابعة في هذا الكتاب، ومرة أخرى إيقاع، دوره هذه المرة نحو يوم. المولّد الرئيسي مجموعة صغيرة من عشرات آلاف العصبونات في قاعدة الدماغ، يزامن إيقاعُها بقية خلايا الجسم.

الدور الذاتي لهذا المولّد عند الإنسان ليس يومًا بالضبط، بل $24.18$ ساعة في المتوسط~\cite{Czeisler1999}. ويضبطه كل يوم الضوءُ الآتي من مستشعرات العين (الفصل~\ref{sec:sensors}). وهذا عند مهندس الإلكترونيات \emph{حلقة مقفلة الطور}: مولّد بتردده الذاتي $\omega_0$ يتزامن مع إشارة خارجية ترددها $\omega$، وفرق الطور $\varphi$ يخضع للمعادلة
\begin{empheq}[box=\keybox]{equation}
\frac{\dd\varphi}{\dd t} \;=\; (\omega-\omega_0) \;-\; K_\varphi\,\sin\varphi .
\label{eq:pll}
\end{empheq}
يبقى التزامن قائمًا إذا كان $|\omega-\omega_0|<K_\varphi$. والمولّد الذي دوره $24.18$ ساعة يحتاج إلى أن ينزاح كل يوم نحو إحدى عشرة دقيقة؛ ويكفي ضوء الصباح عادةً لهذه الإزاحة. أما في الليل القطبي أو في كهف بلا ضوء فلا تزامن، وينجرف الإيقاع إلى دوره الذاتي.

المولّد اليومي ليس مولّد نبضات الساعة في الحاسوب. فهو لا يعدّ خطوات الحساب ولا يزامن نبضات العصبونات. إنه يبدّل \emph{الأوضاع}: النوم واليقظة، ومستويات الهرمونات، وحرارة الجسم، والاستعداد للطعام. إنه سجل بث بطيء من جنس سجلات الفصول~\babelsublr{\ref{sec:serocomputer}--\ref{sec:acet}}، غير أن قيمته تتغير وفق جدول مضبوط على الشمس.

\begin{keyblock}
\begin{proposition}[المخرج الكيميائي]
\label{prop:chemoutput}
المخرج البطيء للآلة مبني من القطع نفسها التي بُني منها المخرج السريع. تتلقى الأعضاء سلكين متعاكسي الإشارة مختلفي السرعة، المسرِّع \nt{NORA} والمكبح \nt{ACET}، ويتلقى الجسم كله معًا إشارات بث عبر الدم، ومنها الأدرينالين \nt{ADRE}. وتثبّت مستوياتِ الهرمونات سلاسلُ مراحل بتغذية راجعة سالبة، دقتها محدودة بالاستقرار؛ وبعض الهرمونات مرمَّز بتردد الدفعات؛ والإيقاع اليومي مولّد بطيء بحلقة مقفلة الطور يضبطها الضوء. بنية الطرق والتجارب على الدفعات وعلى الدور مُقاسة؛ أما تفسير نبضات السلسلة بالتغذية الراجعة نفسها ففرضية.
\end{proposition}
\end{keyblock}

\section{الخلاصة}

\begin{table}[t]
\centering
\footnotesize
\setlength{\tabcolsep}{4pt}
\renewcommand{\arraystretch}{1.15}
\begin{tabular}{>{\raggedright}p{3.1cm}>{\raggedright}p{4.8cm}>{\raggedright\arraybackslash}p{4.9cm}}
\toprule
ماذا يفعل المخرج الكيميائي & كيف & ما المعروف \\
\midrule
يضبط الأعضاء & المسرِّع \nt{NORA} والمكبح \nt{ACET}~\eqref{eq:gasbrake} & مُقاس؛ المكبح أسرع من المسرِّع \\
ينقل الجسم كله إلى وضع الركض & الأدرينالين \nt{ADRE} في الدم & مُقاس \\
يثبّت مستويات الهرمونات & سلسلة مراحل بتغذية راجعة، $K<8$~\eqref{eq:cascadestable} & السلاسل والتغذية الراجعة مُقاسة؛ نشوء النبضات من الحلقة نفسها فرضية \\
ينقل بتردد الدفعات & المستقبل المُتعَب مرشحَ تمرير عالٍ & الاعتماد على الدفعات مُقاس \\
يبدّل الأوضاع على مدار اليوم & مولّد بحلقة مقفلة الطور يضبطها الضوء~\eqref{eq:pll} & الدور والضبط بالضوء مُقاسان \\
\bottomrule
\end{tabular}
\caption{كيف تتحكم الآلة في كيمياء الجسم.}
\label{tab:chemoutput}
\end{table}

للآلة مخرجان (الجدول~\ref{tab:chemoutput}). السريع هو العضلات: أجزاء من الألف من الثانية، معنوَن، إلى كل مفصل. والبطيء هو كيمياء الجسم: ثوانٍ ودقائق وأيام، عبر سلكين إلى الأعضاء وعبر الدم إلى الجميع معًا. والقطع في المخرجين واحدة: سلكان متعاكسا الإشارة، وناظمات إيقاع، وتغذية راجعة وتأخرها، وهي نفسها القطع التي جُمعت منها الآلة من الداخل.

\section{تمارين}

\begin{exercise}[\lvlA]
\label{ex:16:1}
\emph{الدم مرشحًا.} يُطرح هرمون من الدم بعمر نصف $t_{1/2}=10$ دقائق ويُطلق دفعاتٍ قصيرة كل $T=60$ دقيقة. كم ضعفًا تزيد القيمة العظمى للتركيز على الصغرى، وكم ضعفًا تُوهَّن التوافقية الأساسية للدفعات؟ عند أي $t_{1/2}$ تزيد العظمى على الصغرى ضعفين فقط؟
\end{exercise}

\begin{exercise}[\lvlA]
\label{ex:16:2}
\emph{التزامن اليومي.} الدور الذاتي للمولّد~\eqref{eq:pll} يساوي $24.18$ ساعة، والضوء يزيح الطور ساعة في اليوم على الأكثر: $K_\varphi=2\pi/24$ راديان/يوم. أوجد فرق الطور المستقر $\varphi^*$ بالدقائق وزمن الاسترخاء. بعد كم يومًا يصبح عدم التطابق أقل من ساعة إثر قفزة في الإشارة الخارجية مقدارها $\pm6$ ساعات؟
\end{exercise}

\begin{exercise}[\lvlB]
\label{ex:16:3}
\emph{الدقة مقابل الاستقرار.} مُدّدت السلسلة الخطية~\eqref{eq:cascade} إلى $n$ مرحلة متماثلة. أوجد الكسب الحرج $K_c(n)$ وأصغر خطأ ممكن $1/(1+K_c)$ عند $n=3$ و$n=6$. إلامَ يؤول هذا الخطأ عندما $n\to\infty$؟
\end{exercise}

\begin{exercise}[\lvlB]
\label{ex:16:4}
\emph{المسرِّع والمكبح.} في~\eqref{eq:gasbrake} مساهمتا المسرِّع والمكبح مخرجا دارتي RC بـ$\tau_S=2$ ثانية و$\tau_P=0.4$ ثانية، والأمران لا يتجاوزان $80$ و$60$ ضربة في الدقيقة. الإيقاع $f_0=100$؛ في الراحة المكبح $35$ والمسرِّع $0$ و$f=65$. في كم من الزمن يمكن بلوغ $f=120$؟ وماذا لو لم يُرفع المكبح وعمل المسرِّع وحده؟
\end{exercise}

\begin{exercise}[\lvlB]
\label{ex:16:5}
\emph{نمذجة: سلسلة بتأخير.} أضف تأخيرًا $d$ إلى التغذية الراجعة في الدالة \texttt{cascade} من \texttt{models/\allowbreak chem\_models.py} (انسخها ولا تشغّل الملف): ترى المرحلة الأولى $x_3(t-d)$. أوجد $K_c$ والدور عند الحد عند $d/\tau=0$ و$0.5$ و$1$ و$2$، وتحقق بالمحاكاة عند $0.9K_c$ و$1.1K_c$.
\end{exercise}

\begin{exercise}[\lvlB]
\label{ex:16:6}
\emph{متلقٍّ يقرأ الدفعات.} المستقبل يتعب: $y=[c-a]_+$، $\tau_a\,\dd a/\dd t=c-a$، $\tau_a=30$ دقيقة. يصل الهرمون دفعاتٍ مدة كل منها $6$ دقائق بدور $T$ مع الجرعة المتوسطة نفسها $\bar c$. أوجد متوسط الاستجابة عند $T=15$ و$60$ و$120$ دقيقة وعندما $T\to\infty$. وكم يكون عند تركيز ثابت $\bar c$؟
\end{exercise}

\begin{exercise}[\lvlC]
\label{ex:16:7}
\emph{مولّد أم تغذية راجعة؟} اقترح تجربة تميّز النبضات التي تولّدها التغذية الراجعة المتأخرة في السلسلة~\eqref{eq:cascade} من ناظمة إيقاع مستقلة. هل يمكن التمييز بين الفرضيتين بإزاحة الدور إذا لم يمكن تغيير $K$ إلا مرة ونصفًا، وكان الدور يُقاس بدقة $5\%$؟
\end{exercise}
